% !TEX TS-program = lualatex
% !TEX encoding = UTF-8

\documentclass[VesperaleOP.tex]{subfiles}

\ifcsname preamble@file\endcsname
  \setcounter{page}{\getpagerefnumber{M-vop27_commune_martyrum}}
\fi

\begin{document}

\festum{}{Commune unius Martyris}{extra tempus paschale}{Commune unius Martyris extra t. p.}{1}
\addcontentsline{toc}{section}{Commune unius Martyris extra t. p.}

\lineaparva

\titulumrubrica{Ad II Vesperas}
\cantus[Ant. I b]{CUMEXVA}{Super Ps.}
\begin{psalmodia}
\psalmusinpsalmodia{109}
\psalmusinpsalmodia{111}
\psalmusinpsalmodia{112}
\psalmusinpsalmodia{115}
\psalmusinpsalmodia{125}
\end{psalmodia}

\capitulum{Iac. 1}{Beátus vir qui suffert tentatiónem,\pscross{}quóniam, cum probátus fúerit, accípiet corónam vitæ,\psstar{} quam repromísit Deus diligéntibus se.}

\rubrica{In festo II classis:}
\cantus{CUMEXVH1}{Hymnus}

\rubrica{In festis III classis ubi psalmi Dominici dicuntur ad Laudes:}
\cantus{CUMEXVH2}{Hymnus}

\rubrica{In festis III classis ubi psalmi feriales dicuntur ad Laudes:}
\cantus{CUMEXVH3}{Hymnus}

\versiculus{Glória et honóre coronásti eum, Dómine.}{Et constituísti eum super ópera mánuum tuárum.}

\cantus{CUMEXVAM}{Ad Magn.}

\end{document}
